\chapter{2014年四川卷（理）}

\section{单选题}

\begin{problem}
已知集合 \(A = \{x \mid x^2 - x - 2 \leqslant 0\}\)，集合 \(B\) 为整数集，则 \(A \cap B =\)（\quad）

\choices
    {\( \{-1,0,1,2\} \)}
    {\(\{-2, -1, 0, 1\}\)}
    {\(\{0, 1\}\)}
    {\(\{-1,0\}\)}

\answer{A}

\begin{solution}
因
\[
x^2-x-2=(x-2)(x+1)
\]
且二次项系数为正，所以不等式
\[
(x-2)(x+1)\leqslant0
\]
的解集为 \(-1\leqslant x\leqslant2\)，即 \(A=[-1,2]\). 又 \(B=\Z\)，因此
\[
A\cap B=\{-1,0,1,2\}
\]
故选 A.
\end{solution}
\end{problem}

\begin{problem}
在 \(x(1 + x)^6\) 的展开式中，含 \(x^3\) 项的系数为（\quad）

\choices
    {\(30\)}
    {\(20\)}
    {\(15\)}
    {\(10\)}

\answer{C}

\begin{solution}
二项式展开式的一般项为
\[
x(1+x)^6=x\sum_{k=0}^{6}\mathrm C_6^k x^k
=\sum_{k=0}^{6}\mathrm C_6^k x^{k+1}
\]
要得到 \(x^3\) 项，必须有 \(k+1=3\)，即 \(k=2\). 因此所求系数为
\[
\mathrm C_6^2=\frac{6\cdot5}{2\cdot1}=15
\]
故选 C.
\end{solution}
\end{problem}

\begin{problem}
为了得到函数 \( y = \sin (2x + 1) \) 的图象，只需把函数 \( y = \sin 2x \) 的图象上所有的点（\quad）

\choices
    {向左平行移动 \(\frac{1}{2}\) 个单位长度}
    {向右平行移动 \(\frac{1}{2}\) 个单位长度}
    {向左平行移动 1 个单位长度}
    {向右平行移动 1 个单位长度}

\answer{A}

\begin{solution}
将 \(y=\sin2x\) 的图象向左平移 \(h\) 个单位长度后，所得函数为
\[
y=\sin\bigl(2(x+h)\bigr)
\]
令 \(h=\frac12\)，则
\[
\sin\bigl(2(x+\tfrac12)\bigr)=\sin(2x+1)
\]
所以应向左平移 \(\frac12\) 个单位长度，选 A.
\end{solution}
\end{problem}

\begin{problem}
若 \(a > b > 0\)，\(c < d < 0\)，则一定有（\quad）

\choices
    {\(\frac{a}{c} >\frac{b}{d}\)}
    {\(\frac{a}{c} < \frac{b}{d}\)}
    {\(\frac{a}{d} > \frac{b}{c}\)}
    {\(\frac{a}{d} < \frac{b}{c}\)}

\answer{D}

\begin{solution}
令 \(u=-c\)，\(v=-d\)，则 \(u>v>0\). 由 \(a>b>0\)，有
\[
au-bv=a(u-v)+v(a-b)>0
\]
而
\[
\frac{a}{d}-\frac{b}{c}=-\frac{a}{v}+\frac{b}{u}
=\frac{bv-au}{uv}<0
\]
所以
\[
\frac{a}{d}<\frac{b}{c}
\]
选 D.
\end{solution}
\end{problem}

\begin{problem}
执行如图所示的程序框图. 如果输入的 \(x\)，\(y \in \R\)，则输出的 \(S\) 的最大值（\quad）

\texfigure[max width=0.42\linewidth]{img_repaint/2014/sichuan_science/q5_fig1.tex}

\choices
    {\(0\)}
    {\(1\)}
    {\(2\)}
    {\(3\)}

\answer{C}

\begin{solution}
当程序判断条件成立时，
\(x\geqslant0\)，\(y\geqslant0\)，\(x+y\leqslant1\)
并输出 \(S=2x+y\). 在该区域内，
\[
2x+y=x+(x+y)\leqslant x+1\leqslant2
\]
且 \(x=1\)，\(y=0\) 时取到 \(2\)，所以这一分支的最大值为 \(2\). 当判断条件不成立时，程序直接输出 \(S=1\)，小于 \(2\). 因此输出的 \(S\) 的最大值为 \(2\)，选 C.
\end{solution}
\end{problem}

\begin{problem}
六个人从左至右排成一行，最左端只能排甲或乙，最右端不能排甲，则不同的排法共有（\quad）

\choices
    {\(192\text{ 种}\)}
    {\(216\text{ 种}\)}
    {\(240\text{ 种}\)}
    {\(288\text{ 种}\)}

\answer{B}

\begin{solution}
按最左端所排人员分类计数.

若最左端排甲，则甲已经不在最右端，剩余五人可以任意排列，排法数为
\[
5!=120
\]
若最左端排乙，则剩余五人中除甲外还有四人可以排在最右端，最右端有 \(4\) 种选法；确定最右端后，其余四人任意排列，有 \(4!\) 种，故排法数为
\[
4\cdot4!=96
\]
两类互不相交，故总排法数为
\[
5!+4\cdot4!=120+96=216
\]
选 B.
\end{solution}
\end{problem}

\begin{problem}
平面向量 \(\bs{a} = (1, 2)\)，\(\bs{b} = (4, 2)\)，\(\bs{c} = m\bs{a} + \bs{b}\)（\(m \in \R\)），且 \(\bs{c}\) 与 \(\bs{a}\) 的夹角等于 \(\bs{c}\) 与 \(\bs{b}\) 的夹角，则 \(m =\)（\quad）

\choices
    {\(-2\)}
    {\(-1\)}
    {\(1\)}
    {\(2\)}

\answer{D}

\begin{solution}
由
\[
\bs c=m\bs a+\bs b=(m+4,2m+2)
\]
向量 \(\bs c\) 不可能为零，因为 \(m+4=0\) 与 \(2m+2=0\) 不可能同时成立. 又
\(|\bs a|=\sqrt5\)，\(|\bs b|=\sqrt{20}=2\sqrt5\).
两夹角相等，所以它们的余弦相等：
\[
\frac{\bs c\cdot\bs a}{|\bs c||\bs a|}=\frac{\bs c\cdot\bs b}{|\bs c||\bs b|}
\]
计算两个数量积：
\[
\bs c\cdot\bs a=(m+4)+2(2m+2)=5m+8
\]
\[
\bs c\cdot\bs b=4(m+4)+2(2m+2)=8m+20
\]
代入并约去正数 \(|\bs c|\sqrt5\)，得到
\[
\frac{5m+8}{\sqrt5}=\frac{8m+20}{2\sqrt5}
\]
即
\(2(5m+8)=8m+20\)，\(m=2\).
选 D.
\end{solution}
\end{problem}

\begin{problem}
如图，在正方体 \(ABCD - A_{1}B_{1}C_{1}D_{1}\) 中，点 \(O\) 为线段 \(BD\) 的中点. 设点 \(P\) 在线段 \(CC_{1}\) 上，直线 \(OP\) 与平面 \(A_{1}BD\) 所成的角为 \(\alpha\)，则 \(\sin \alpha\) 的取值范围是（\quad）

\bitmapfigure[max width=0.42\linewidth]{img/2014/sichuan_science/q8_fig1.png}

\choices
    {\( \left[\frac{\sqrt{3}}{3},1\right] \)}
    {\( \left[\frac{\sqrt{6}}{3},1\right] \)}
    {\( \left[\frac{\sqrt{6}}{3},\frac{2\sqrt{2}}{3}\right] \)}
    {\( \left[\frac{2\sqrt{2}}{3},1\right] \)}

\answer{B}

\begin{solution}
设正方体棱长为 \(1\)，建立空间直角坐标系：
\(A=(0,0,0)\)，\(B=(1,0,0)\)，\(D=(0,1,0)\)，\(A_1=(0,0,1)\).
于是
\(C=(1,1,0)\)，\(O=(\tfrac12,\tfrac12,0)\).
设 \(P=(1,1,t)\)，则 \(0\leqslant t\leqslant1\).

平面 \(A_1BD\) 内有两个方向向量
\(\overrightarrow{A_1B}=(1,0,-1)\)，\(\overrightarrow{A_1D}=(0,1,-1)\)
故可取法向量
\[
\bs n=\overrightarrow{A_1B}\times\overrightarrow{A_1D}=(1,1,1)
\]
又
\[
\overrightarrow{OP}=P-O=(\tfrac12,\tfrac12,t)
\]
直线与平面所成角的正弦等于方向向量与平面法向量夹角的余弦，因此
\[
\sin\alpha=\frac{|\overrightarrow{OP}\cdot\bs n|}{|\overrightarrow{OP}|\,|\bs n|}=\frac{1+t}{\sqrt{3(t^2+\frac12)}}
\]
令
\(h(t)=\frac{(1+t)^2}{3(t^2+\frac12)}\)，\(h'(t)=\frac{2(1+t)(\frac12-t)}{3(t^2+\frac12)^2}\).
由于 \(1+t>0\)，所以 \(h(t)\) 在 \([0,\frac12]\) 上递增，在 \([\frac12,1]\) 上递减，最大值为
\[
h(\tfrac12)=1
\]
两端点处
\(\sin\alpha\big|_{t=0}=\frac{\sqrt6}{3}\)，\(\sin\alpha\big|_{t=1}=\frac{2\sqrt2}{3}>\frac{\sqrt6}{3}\).
所以最小值在 \(t=0\) 处取得，且
\[
\sin\alpha\in\left[\frac{\sqrt6}{3},1\right]
\]
选 B.
\end{solution}
\end{problem}

\begin{problem}
已知 \(f(x) = \ln (1 + x) - \ln (1 - x)\)，\(x \in (-1, 1)\). 现有下列命题：

\begin{circlelist}
\item \( f(-x) = -f(x) \)；
\item \( f\left(\frac{2x}{x^{2}+1}\right) = 2f(x) \)；
\item \( |f(x)| \geqslant 2|x| \).
\end{circlelist}

其中所有正确命题的序号是（\quad）

\choices
    {\circled{1}\circled{2}\circled{3}}
    {\circled{2}\circled{3}}
    {\circled{1}\circled{3}}
    {\circled{1}\circled{2}}

\answer{A}

\begin{solution}
先验证定义域：当 \(-1<x<1\) 时，\(1+x>0\)，\(1-x>0\)，故下列变形均有意义.

对第一个命题，
\[
f(-x)=\ln(1-x)-\ln(1+x)
=-\bigl(\ln(1+x)-\ln(1-x)\bigr)=-f(x)
\]
所以第一个命题正确.

对第二个命题，令 \(u=\frac{2x}{1+x^2}\). 因为
\[
1-u^2=1-\frac{4x^2}{(1+x^2)^2}
=\frac{(1-x^2)^2}{(1+x^2)^2}>0
\]
所以 \(|u|<1\)，可以代入 \(f\). 于是
\[
f(u)=\ln\frac{1+u}{1-u}=\ln\frac{(1+x)^2}{(1-x)^2}=2\ln\frac{1+x}{1-x}=2f(x)
\]
所以第二个命题正确.

对第三个命题，先设 \(0\leqslant x<1\). 有
\(f'(x)=\frac1{1+x}+\frac1{1-x}=\frac2{1-x^2}\geqslant2\)，\(f(0)=0\).
因此
\[
f(x)-f(0)=\int_0^x f'(t)\,\mathrm{d}t\geqslant\int_0^x2\,\mathrm{d}t=2x
\]
当 \(-1<x<0\) 时，由奇函数性质 \(f(x)=-f(-x)\)，且 \(-x\in(0,1)\)，从而
\[
|f(x)|=|f(-x)|\geqslant2|-x|=2|x|
\]
在 \(x=0\) 时等号成立，故第三个命题也正确. 三条命题均正确，选 A.
\end{solution}
\end{problem}

\begin{problem}
已知 \(F\) 是抛物线 \(y^{2}=x\) 的焦点，点 \(A\)，\(B\) 在该抛物线上且位于 \(x\) 轴的两侧，\(\overrightarrow{OA} \cdot \overrightarrow{OB}=2\)（其中 \(O\) 为坐标原点），则 \(\triangle ABO\) 与 \(\triangle AFO\) 面积之和的最小值是（\quad）

\choices
    {\(2\)}
    {\(3\)}
    {\(\frac{17\sqrt{2}}{8}\)}
    {\(\sqrt{10}\)}

\answer{B}

\begin{solution}
抛物线 \(y^2=x\) 可参数表示为 \((t^2,t)\). 由于 \(A\)，\(B\) 位于 \(x\) 轴两侧，设
\(A=(p^2,p)\)，\(B=(q^2,q)\)，\(p>0\)，\(\ q<0\).
由数量积条件
\[
\overrightarrow{OA}\cdot\overrightarrow{OB}=p^2q^2+pq=pq(pq+1)=2
\]
令 \(r=pq\)，则 \(r(r+1)=2\)，即
\[
(r-1)(r+2)=0
\]
因为 \(p>0\)，\(q<0\)，所以 \(r=pq<0\)，只能取 \(pq=-2\). 抛物线 \(y^2=x\) 的焦点为
\[
F=(\tfrac14,0)
\]
三角形 \(ABO\) 的面积为
\[
S_{\triangle ABO}=\frac12|p^2q-pq^2|
=\frac12|pq(p-q)|=p-q
\]
其中最后一步用了 \(pq=-2\) 以及 \(p-q>0\). 又
\[
S_{\triangle AFO}=\frac12\left|\det\begin{pmatrix}\frac14&0\\p^2&p\end{pmatrix}\right|=\frac p8
\]
由 \(q=-\frac2p\)，面积和为
\[
S=p-q+\frac p8=\frac98p+\frac2p
\]
由基本不等式
\[
S\geqslant2\sqrt{\frac98p\cdot\frac2p}=3
\]
当且仅当
\[
\frac98p=\frac2p
\]
即 \(p=\frac43\)（此时 \(q=-\frac32\)）时取等号. 故最小值为 \(3\)，选 B.
\end{solution}
\end{problem}

\section{填空题}

\begin{problem}
复数 \( \frac{2-2i}{1+i}= \)\fillinblank{}.

\answer{\(-2i\)}

\begin{solution}
分母 \(1+i\) 的共轭复数为 \(1-i\). 分子、分母同乘 \(1-i\)，得
\[
\frac{2-2i}{1+i}=\frac{(2-2i)(1-i)}{(1+i)(1-i)}=\frac{-4i}{2}=-2i
\]
因此填 \(-2i\).
\end{solution}
\end{problem}

\begin{problem}
设 \( f(x) \) 是定义在 \(\R\) 上的周期为 2 的函数，当 \( x \in [-1, 1) \) 时， \( f(x) = \begin{cases}
-4x^{2} + 2, & -1 \leqslant x < 0,\\
x, & 0 \leqslant x < 1,
\end{cases} \) 则 \( f\left(\frac{3}{2}\right) = \)\fillinblank{}.

\answer{\(1\)}

\begin{solution}
周期为 \(2\) 表示对任意 \(x\)，都有 \(f(x+2)=f(x)\). 由于
\[
\frac32-2=-\frac12\in[-1,1)
\]
所以
\[
f\left(\frac32\right)=f\left(-\frac12\right)
=-4\left(-\frac12\right)^2+2=-1+2=1
\]
因此填 \(1\).
\end{solution}
\end{problem}

\begin{problem}
如图，从气球 \(A\) 上测得正前方的河流的两岸 \(B\)，\(C\) 的俯角分别为 \(67^{\circ}\)，\(30^{\circ}\)，此时气球的高是 \(46\mathrm{~m}\)，则河流的宽度 \(BC\) 约等于\fillinblank{}\(\mathrm{m}\). （用四舍五入法将结果精确到个位. 参考数据：\(\sin 67^{\circ} \approx 0.92\)，\(\cos 67^{\circ} \approx 0.39\)，\(\sin 37^{\circ} \approx 0.60\)，\(\cos 37^{\circ} \approx 0.80\)，\(\sqrt{3} \approx 1.73\)）

\bitmapfigure[max width=0.42\linewidth]{img/2014/sichuan_science/q13_fig1.png}

\answer{\(60\)}

\begin{solution}
设气球 \(A\) 在河岸所在水平面上的正投影为 \(H\). 由图中 \(67^\circ>30^\circ\)，可知 \(B\) 比 \(C\) 更靠近 \(H\)，所以
\[
BC=HC-HB
\]
俯角等于相应的仰角，在直角三角形 \(AHB\) 与 \(AHC\) 中，
\(HB=46\cot67^{\circ}\)，\(HC=46\cot30^{\circ}=46\sqrt3\).
所以
\[
BC=HC-HB=46\left(\sqrt3-\frac{\cos67^{\circ}}{\sin67^{\circ}}\right)
\approx46\left(1.73-\frac{0.39}{0.92}\right)\approx60.1
\]
四舍五入到个位得
\[
BC\approx60\mathrm{~m}
\]
\end{solution}
\end{problem}

\begin{problem}
设 \(m \in \R\)，过定点 \(A\) 的动直线 \(x + my = 0\) 和过定点 \(B\) 的动直线 \(mx - y - m + 3 = 0\) 交于点 \(P(x, y)\)，则 \(|PA| \cdot |PB|\) 的最大值是\fillinblank{}.

\answer{\(5\)}

\begin{solution}
由 \(x+my=0\) 可知第一条直线恒过 \(A=(0,0)\). 第二条直线在 \(x=1\)，\(y=3\) 时恒成立：
\[
m\cdot1-3-m+3=0
\]
故它恒过 \(B=(1,3)\). 交点 \(P(x,y)\) 满足
\(x+my=0\)，\(mx-y-m+3=0\)
解得
\(x=\frac{m(m-3)}{m^2+1}\)，\(y=\frac{3-m}{m^2+1}\).
因为 \(m^2+1>0\)，由此
\[
|PA|=\sqrt{x^2+y^2}=\frac{|m-3|}{\sqrt{m^2+1}}
\]
又因 \(P\)，\(B\) 在第二条直线上，直接计算得
\(x-1=-\frac{3m+1}{m^2+1}\)，\(y-3=m(x-1)\)
从而
\[
|PB|=\sqrt{(x-1)^2+(y-3)^2}=\frac{|3m+1|}{\sqrt{m^2+1}}
\]
于是
\[
|PA||PB|=\frac{|3m^2-8m-3|}{m^2+1}
\]
为证明其不超过 \(5\)，计算
\[
25(m^2+1)^2-(3m^2-8m-3)^2=4(m+2)^2(2m-1)^2\geqslant0
\]
由于 \(m^2+1>0\)，上式等价于
\[
|3m^2-8m-3|\leqslant5(m^2+1)
\]
故 \(|PA||PB|\leqslant5\). 当 \(m=-2\) 或 \(m=\frac12\) 时等号成立，所以最大值为 \(5\).
\end{solution}
\end{problem}

\begin{problem}
以 \(A\) 表示值域为 \(\R\) 的函数组成的集合，\(B\) 表示具有如下性质的函数 \(\varphi(x)\) 组成的集合：对于函数 \(\varphi(x)\)，存在一个正数 \(M\)，使得函数 \(\varphi(x)\) 的值域包含于区间 \([-M, M]\). 例如，当 \(\varphi_1(x) = x^3\)，\(\varphi_2(x) = \sin x\) 时，\(\varphi_1(x) \in A\)，\(\varphi_2(x) \in B\). 现有如下命题：

\begin{circlelist}
\item 设函数 \(f(x)\) 的定义域为 \(D\)，则“\(f(x) \in A\)”的充要条件是“\(\forall b \in \R\)，\(\exists a \in D\)，\(f(a)=b\)”；
\item 函数 \(f(x) \in B\) 的充要条件是 \(f(x)\) 有最大值和最小值；
\item 若函数 \(f(x)\)，\(g(x)\) 的定义域相同，且 \(f(x) \in A\)，\(g(x) \in B\)，则 \(f(x) + g(x) \notin B\)；
\item 若函数 \(f(x) = a\ln(x + 2) + \frac{x}{x^{2} + 1}\)（\(x > -2\)，\(a \in \R\)）有最大值，则 \(f(x) \in B\).
\end{circlelist}

其中的真命题有\fillinblank{}.（写出所有真命题的序号）

\answer{\(\circled{1}\circled{3}\circled{4}\)}

\begin{solution}
第一个命题正是“函数值域为 \(\R\)”的定义：值域为 \(\R\) 等价于任意 \(b\in\R\) 都能在定义域 \(D\) 中找到 \(a\)，使 \(f(a)=b\). 故第一个命题正确.

第二个命题不正确. 若函数有最大值和最小值，当然可以推出值域有界；但有界函数不一定取得最大值和最小值. 例如 \(f(x)=x\) 在定义域 \((0,1)\) 上的值域为 \((0,1)\)，它被 \([-1,1]\) 包含，却没有最大值和最小值. 所以“有界”不是“有最大值和最小值”的充要条件.

对第三个命题，设 \(g(x)\) 的值域包含于 \([-M,M]\)，则对所有 \(x\)，有 \(-M\leqslant g(x)\leqslant M\). 因为 \(f(x)\) 的值域为 \(\R\)，对任意 \(K>0\)，都存在定义域中的点 \(x_1\)，\(x_2\)，使
\(f(x_1)>K+M\)，\(f(x_2)<-K-M\).
于是
\(f(x_1)+g(x_1)>K\)，\(f(x_2)+g(x_2)<-K\).
因此 \(f+g\) 无界，不属于 \(B\)，第三个命题正确.

对第四个命题，若 \(a>0\)，则当 \(x\to+\infty\) 时，\(a\ln(x+2)\to+\infty\)，且 \(\frac{x}{x^2+1}\to0\)，所以 \(f(x)\to+\infty\)，不可能有最大值. 若 \(a<0\)，则当 \(x\to-2^+\) 时，\(\ln(x+2)\to-\infty\)，从而 \(a\ln(x+2)\to+\infty\)，第二项有有限极限，因此同样不可能有最大值. 所以若 \(f\) 有最大值，只能有 \(a=0\).

此时
\[
f(x)=\frac{x}{x^2+1}
\]
由
\[
x^2+1\geqslant2|x|
\]
可得
\[
\left|\frac{x}{x^2+1}\right|\leqslant\frac12
\]
所以此时值域包含于 \([-\frac12,\frac12]\)，即 \(f(x)\in B\). 因此第四个命题正确，真命题的序号为 1，3，4.
\end{solution}
\end{problem}

\section{解答题}

\begin{problem}
已知函数 \( f(x) = \sin \left(3x + \frac{\pi}{4}\right) \).
\begin{enumerate}
\item 求 \( f(x) \) 的单调递增区间；
\item 若 \(\alpha\) 是第二象限角，\(f\left(\frac{\alpha}{3}\right) = \frac{4}{5}\cos \left(\alpha + \frac{\pi}{4}\right)\cos 2\alpha\)，求 \(\cos \alpha - \sin \alpha\) 的值.
\end{enumerate}
\end{problem}

\begin{solution}
\begin{enumerate}
\item 函数 \(y=\sin t\) 在每个区间
\(\left[-\frac\pi2+2k\pi,\frac\pi2+2k\pi\right]\)，\(k\in\Z\)
上单调递增. 令 \(t=3x+\frac\pi4\)，由于 \(3>0\)，只需解
\[
-\frac\pi2+2k\pi\leqslant3x+\frac\pi4\leqslant\frac\pi2+2k\pi
\]
解得
\(x\in\left[\frac{2k\pi}{3}-\frac\pi4,\frac{2k\pi}{3}+\frac\pi{12}\right]\)，\(k\in\Z\).
所以 \(f(x)\) 的单调递增区间为
\(\left[\frac{2k\pi}{3}-\frac\pi4,\frac{2k\pi}{3}+\frac\pi{12}\right]\)，\(k\in\Z\).

\item 由
\[
f\left(\frac\alpha3\right)=\sin\left(3\cdot\frac\alpha3+\frac\pi4\right)
=\sin\left(\alpha+\frac\pi4\right)
\]
已知条件化为
\[
\sin\left(\alpha+\frac\pi4\right)=\frac45\cos\left(\alpha+\frac\pi4\right)\cos2\alpha
\]
令
\(u=\cos\alpha-\sin\alpha\)，\(v=\cos\alpha+\sin\alpha\).
则
\(\cos\left(\alpha+\frac\pi4\right)=\frac{u}{\sqrt2}\)，\(\sin\left(\alpha+\frac\pi4\right)=\frac{v}{\sqrt2}\)，\(\cos2\alpha=uv\).
代入上式并约去 \(\sqrt2\)，得到
\[
v=\frac45u^2v
\]
分情况讨论：若 \(v=0\)，则 \(\cos\alpha+\sin\alpha=0\). 因 \(\alpha\) 为第二象限角，故 \(\alpha=\frac{3\pi}{4}\)，从而
\[
u=\cos\frac{3\pi}{4}-\sin\frac{3\pi}{4}=-\sqrt2
\]
若 \(v\ne0\)，则可约去 \(v\)，得到 \(u^2=\frac54\). 又因 \(\alpha\) 为第二象限角，\(\cos\alpha<0\)，\(\sin\alpha>0\)，所以 \(u<0\)，从而
\[
u=-\frac{\sqrt5}{2}
\]
因此
\[
\cos\alpha-\sin\alpha=-\sqrt2\quad\text{或}\quad-\frac{\sqrt5}{2}
\]
\end{enumerate}
\end{solution}

\begin{problem}
一款击鼓小游戏的规则如下：每盘游戏都需要击鼓三次，每次击鼓要么出现一次音乐，要么不出现音乐；每盘游戏击鼓三次后，出现一次音乐获得 10 分，出现两次音乐获得 20 分，出现三次音乐获得 100 分，没有出现音乐则扣除 200 分（即获得 \(-200\) 分）. 设每次击鼓出现音乐的概率为 \(\frac{1}{2}\)，且各次击鼓出现音乐相互独立.
\begin{enumerate}
\item 设每盘游戏获得的分数为 \(X\)，求 \(X\) 的分布列；
\item 玩三盘游戏，至少有一盘出现音乐的概率是多少？
\item 玩过这款游戏的许多人都发现，若干盘游戏后，与最初的分数相比，分数没有增加反而减少了. 请运用概率统计的相关知识分析分数减少的原因.
\end{enumerate}
\end{problem}

\begin{solution}
\begin{enumerate}
\item 设一盘游戏出现音乐的次数为 \(Y\). 由于每次击鼓独立且出现音乐的概率为 \(\frac12\)，所以
\[
Y\sim B\left(3,\frac12\right)
\]
从而
\(P(Y=k)=\mathrm C_3^k\left(\frac12\right)^3\)，\(k=0\)，\(1\)，\(2\)，\(3\).
逐项计算得
\(P(Y=0)=\frac18\)，\(P(Y=1)=\frac38\)，\(P(Y=2)=\frac38\)，\(P(Y=3)=\frac18\).
按得分规则，\(Y=0\)，\(1\)，\(2\)，\(3\) 时分别有 \(X=-200\)，\(10\)，\(20\)，\(100\)，故
\begin{center}
\begin{tabular}{c|cccc}
 \(X\)&\(-200\)&\(10\)&\(20\)&\(100\)\\ \hline
 \(P\)&\(\frac18\)&\(\frac38\)&\(\frac38\)&\(\frac18\)
\end{tabular}
\end{center}

\item 设 \(A_i\) 表示第 \(i\) 盘没有出现音乐. 一盘没有出现音乐即三次击鼓均未出现音乐，因此
\[
P(A_i)=\left(\frac12\right)^3=\frac18
\]
各盘所用的击鼓次数彼此独立，所以 \(A_1\)，\(A_2\)，\(A_3\) 相互独立. 三盘都没有出现音乐的概率为
\[
P(A_1\cap A_2\cap A_3)=\left(\frac18\right)^3=\frac1{512}
\]
所求事件是其对立事件，故至少有一盘出现音乐的概率为
\[
1-\frac1{512}=\frac{511}{512}
\]

\item 由分布列，每盘游戏得分的数学期望为
\[
E(X)=-200\cdot\frac18+10\cdot\frac38+20\cdot\frac38+100\cdot\frac18=-\frac54<0
\]
若玩 \(n\) 盘，设总得分为 \(S_n=X_1+\cdots+X_n\)，则
\[
E(S_n)=E(X_1)+\cdots+E(X_n)=-\frac54n
\]
因此长期平均得分的期望为 \(-\frac54\)，且由大数定律，实际平均得分在盘数较多时趋近于 \(-\frac54\). 虽然出现三次音乐时可得 \(100\) 分，但没有音乐的概率为 \(\frac18\)，一次就扣 \(200\) 分；正得分的概率加权平均不足以抵消这项损失. 所以分数在若干盘后没有增加反而减少，是负期望导致的正常统计现象，短期内的偶然波动不改变这一长期趋势.
\end{enumerate}
\end{solution}

\begin{problem}
三棱锥 \(A - BCD\) 及其侧视图、俯视图如图所示. 设 \(M\)，\(N\) 分别为线段 \(AD\)，\(AB\) 的中点，\(P\) 为线段 \(BC\) 上的点，且 \(MN \perp NP\).
\begin{enumerate}
\item 证明：\(P\) 为线段 \(BC\) 的中点；
\item 求二面角 \(A - NP - M\) 的余弦值.
\end{enumerate}

\bitmapfigure[max width=0.78\linewidth]{img/2014/sichuan_science/q18_fig1.png}
\end{problem}

\begin{solution}
由侧视图可知三角形 \(ABD\) 为边长 \(2\) 的等边三角形，由俯视图可知三角形 \(BCD\) 也为边长 \(2\) 的等边三角形. 取空间直角坐标系，使
\(B=(-1,0,0)\)，\(D=(1,0,0)\)，\(C=(0,\sqrt3,0)\)，\(A=(0,0,\sqrt3)\).
这样 \(BD=BC=CD=AB=AD=2\)，且两个投影视图中的边长关系均满足. 于是
\(M=\left(\frac12,0,\frac{\sqrt3}{2}\right)\)，\(N=\left(-\frac12,0,\frac{\sqrt3}{2}\right)\).

\begin{enumerate}
\item
设 \(P\) 将 \(BC\) 按参数 \(t\) 分割，即
\(P=(-1+t,\sqrt3t,0)\)，\(0\leqslant t\leqslant1\).
因此
\(\overrightarrow{MN} =(-1,0,0)\)，\(\overrightarrow{NP}=\left(t-\frac12,\sqrt3t,-\frac{\sqrt3}{2}\right)\).
由 \(MN\perp NP\) 得
\[
\overrightarrow{MN}\cdot\overrightarrow{NP}
=-\left(t-\frac12\right)=0
\]
所以 \(t=\frac12\)，即 \(P\) 为 \(BC\) 的中点.

\item
此时
\[
P=\left(-\frac12,\frac{\sqrt3}{2},0\right)
\]
取与相应向量同方向的倍数，有
\(\overrightarrow{NP}\parallel(0,1,-1)\)，\(\overrightarrow{NA}\parallel(1,0,\sqrt3)\)，\(\overrightarrow{NM}\parallel(1,0,0)\).
二面角 \(A-NP-M\) 的两个半平面分别是 \(ANP\) 与 \(MNP\). 因此这两个平面的法向量可分别取为
\[
\bs n_1=(0,1,-1)\times(1,0,\sqrt3)=(\sqrt3,-1,-1)
\]
\[
\bs n_2=(1,0,0)\times(0,1,-1)=(0,1,1)
\]
设二面角 \(A-NP-M\) 的大小为 \(\theta\)，则二面角的余弦取两法向量夹角余弦的绝对值：
\[
\cos\theta=\frac{|\bs n_1\cdot\bs n_2|}{|\bs n_1||\bs n_2|}
=\frac{2}{\sqrt5\sqrt2}=\frac{\sqrt{10}}5
\]
故二面角的余弦值为 \(\frac{\sqrt{10}}5\).
\end{enumerate}
\end{solution}

\begin{problem}
设等差数列 \(\{a_{n}\}\) 的公差为 \(d\)，点 \((a_{n}, b_{n})\) 在函数 \(f(x) = 2^{x}\) 的图象上 （\(n \in \N^{*}\)）.
\begin{enumerate}
\item 若 \(a_{1} = -2\)，点 \((a_{8}, 4b_{7})\) 在函数 \(f(x)\) 的图象上，求数列 \(\{a_{n}\}\) 的前 \(n\) 项和 \(S_{n}\)；
\item 若 \(a_{1} = 1\)，函数 \(f(x)\) 的图象在点 \((a_{2}, b_{2})\) 处的切线在 \(x\) 轴上的截距为 \(2 - \frac{1}{\ln 2}\)，求数列 \(\left\{\frac{a_n}{b_n}\right\}\) 的前 \(n\) 项和 \(T_{n}\).
\end{enumerate}
\end{problem}

\begin{solution}
\begin{enumerate}
\item 因为点 \((a_n,b_n)\) 在 \(y=2^x\) 的图象上，所以
\[
b_n=2^{a_n}
\]
由点 \((a_8,4b_7)\) 也在该图象上，得
\[
2^{a_8}=4\cdot2^{a_7}=2^{a_7+2}
\]
指数函数 \(2^x\) 在实数范围内严格递增，所以
\[
a_8=a_7+2
\]
又 \(a_8-a_7=d\)，从而 \(d=2\). 因此
\[
S_n=na_1+\frac{n(n-1)}2d
=-2n+\frac{n(n-1)}2\cdot2
=n^2-3n
\]

\item 函数 \(y=2^x\) 的导数为 \(y'=2^x\ln2\)，所以在 \((a_2,b_2)\) 处的切线斜率为 \(b_2\ln2\)，切线方程为
\[
y-b_2=b_2\ln2\,(x-a_2)
\]
令 \(y=0\)，并用 \(b_2=2^{a_2}\ne0\)，得切线在 \(x\) 轴上的截距
\[
a_2-\frac1{\ln2}=2-\frac1{\ln2}
\]
所以 \(a_2=2\). 由 \(a_1=1\) 得 \(d=a_2-a_1=1\)，故
\(a_n=a_1+(n-1)d=n\)，\(b_n=2^n\).
于是
\[
T_n=\sum_{k=1}^n\frac{k}{2^k}
\]
为计算该和，令
\[
U_n=\sum_{k=1}^n\frac{k}{2^k}
\]
则
\[
2U_n=\sum_{k=1}^n\frac{k}{2^{k-1}}
=1+\sum_{k=2}^n\frac{k}{2^{k-1}}
\]
两式相减得
\[
U_n=1+\sum_{k=2}^n\frac1{2^{k-1}}-\frac{n}{2^n}
=1+\sum_{j=1}^{n-1}\frac1{2^j}-\frac{n}{2^n}
=2-\frac{n+2}{2^n}
\]
所以
\[
T_n=2-\frac{n+2}{2^n}
=\frac{2^{n+1}-n-2}{2^n}
\]
\end{enumerate}
\end{solution}

\begin{problem}
已知椭圆 \(C: \frac{x^2}{a^2} + \frac{y^2}{b^2} = 1\)（\(a > b > 0\)）的焦距为 4，其短轴的两个端点与长轴的一个端点构成正三角形.
\begin{enumerate}
\item 求椭圆 \(C\) 的标准方程；
\item 设 \(F\) 为椭圆 \(C\) 的左焦点，\(T\) 为直线 \(x = -3\) 上任意一点，过 \(F\) 作 \(TF\) 的垂线交椭圆 \(C\) 于点 \(P\)，\(Q\).
\begin{enumerate}
\item 证明：\(OT\) 平分线段 \(PQ\)（其中 \(O\) 为坐标原点）；
\item 当 \(\frac{|TF|}{|PQ|}\) 最小时，求点 \(T\) 的坐标.
\end{enumerate}
\end{enumerate}
\end{problem}

\begin{solution}
\begin{enumerate}
\item 设椭圆焦半距为 \(c\). 由焦距为 \(4\)，得 \(2c=4\)，所以 \(c=2\). 椭圆参数满足
\[
a^2-b^2=c^2=4
\]
短轴端点为 \((0,b)\)，\((0,-b)\)，取长轴右端点 \((a,0)\). 短轴两个端点间距离为 \(2b\)，而任一短轴端点到 \((a,0)\) 的距离为 \(\sqrt{a^2+b^2}\). 正三角形条件给出
\(\sqrt{a^2+b^2}=2b\)，\(a^2=3b^2\).
联立 \(a^2-b^2=4\) 与 \(a^2=3b^2\)，得
\(2b^2=4\)，\(b^2=2\)，\(a^2=6\).
所以椭圆的标准方程为
\[
\frac{x^2}{6}+\frac{y^2}{2}=1
\]

\item
\begin{enumerate}
\item 由（1）知左焦点 \(F=(-2,0)\). 设 \(T=(-3,t)\)，其中 \(t\in\R\). 则 \(\overrightarrow{TF}=(1,-t)\)，与之垂直的一个方向向量为 \((t,1)\)，故过 \(F\) 且垂直于 \(TF\) 的直线可参数表示为
\[
(x,y)=(-2,0)+\lambda(t,1)=(-2+t\lambda,\lambda)
\]
将其代入椭圆方程
\[
\frac{x^2}{6}+\frac{y^2}{2}=1
\]
得
\[
(t^2+3)\lambda^2-4t\lambda-2=0
\]
由于常数项为 \(-2\)，该方程有两个异号实根，正好对应弦端点 \(P\)，\(Q\). 设两根为 \(\lambda_1\)，\(\lambda_2\)，由韦达定理
\[
\lambda_1+\lambda_2=\frac{4t}{t^2+3}
\]
弦中点对应参数 \(\lambda_0=\frac{\lambda_1+\lambda_2}{2}=\frac{2t}{t^2+3}\)，所以弦 \(PQ\) 的中点为
\[
N=\left(-2+\frac{2t^2}{t^2+3},\frac{2t}{t^2+3}\right)=\left(-\frac6{t^2+3},\frac{2t}{t^2+3}\right)
\]
由上式有 \(y=-\frac t3x\)，而 \(O=(0,0)\)，\(T=(-3,t)\) 确定的直线 \(OT\) 也满足 \(y=-\frac t3x\). 所以 \(N\in OT\)，即 \(OT\) 平分线段 \(PQ\).

\item 由上面的二次方程，根的差满足
\[
|\lambda_1-\lambda_2|
=\frac{\sqrt{(-4t)^2-4(t^2+3)(-2)}}{t^2+3}
=\frac{2\sqrt6\sqrt{t^2+1}}{t^2+3}
\]
参数方向向量 \((t,1)\) 的长度为 \(\sqrt{t^2+1}\)，因此
\(|PQ|=\frac{2\sqrt6(t^2+1)}{t^2+3}\)，\(|TF|=\sqrt{t^2+1}\).
于是
\[
\left(\frac{|TF|}{|PQ|}\right)^2=\frac{(t^2+3)^2}{24(t^2+1)}
\]
令 \(u=t^2\geqslant0\)，则需使
\[
\Phi(u)=\frac{(u+3)^2}{24(u+1)}
\]
最小. 求导得
\[
\Phi'(u)=\frac{(u+3)(u-1)}{24(u+1)^2}
\]
在 \(u\geqslant0\) 上，\(\Phi'(u)<0\)（\(0\leqslant u<1\)），\(\Phi'(1)=0\)，\(\Phi'(u)>0\)（\(u>1\)），所以 \(\Phi\) 在 \(u=1\) 处取得最小值. 因此 \(t^2=1\)，即 \(t=\pm1\). 故点 \(T\) 的坐标为
\[
(-3,1)\quad\text{或}\quad(-3,-1)
\]
\end{enumerate}
\end{enumerate}
\end{solution}

\begin{problem}
已知函数 \( f(x) = \e^{x} - ax^{2} - bx - 1 \)，其中 \(a\)，\(b \in \R\)，\(\e = 2.71828\cdots\) 为自然对数的底数.
\begin{enumerate}
\item 设 \( g(x) \) 是函数 \( f(x) \) 的导函数，求函数 \( g(x) \) 在区间 \( [0,1] \) 上的最小值；
\item 若 \( f(1)=0 \)，函数 \( f(x) \) 在区间 \( (0,1) \) 内有零点，求 a 的取值范围.
\end{enumerate}
\end{problem}

\begin{solution}
\begin{enumerate}
\item 由 \(g(x)=f'(x)\)，得
\(g(x)=\e^x-2ax-b\)，\(g'(x)=\e^x-2a\).
因为 \(\e^x\) 在 \([0,1]\) 上递增，所以按 \(a\) 的取值分类.
\[
\begin{cases}
g'(x)\geqslant0\quad(0\leqslant x\leqslant1),&a\leqslant\frac12,\\
g'(\ln(2a))=0,\quad \ln(2a)\in(0,1),&\frac12<a<\frac\e2,\\
g'(x)\leqslant0\quad(0\leqslant x\leqslant1),&a\geqslant\frac\e2.
\end{cases}
\]
当 \(a\leqslant\frac12\) 时，\(g'(x)\geqslant0\)，\(g\) 在 \([0,1]\) 上递增，最小值为
\[
g(0)=1-b
\]
当 \(\frac12<a<\frac\e2\) 时，\(g'\) 在 \(x=\ln(2a)\) 处由负变正，故最小值为
\[
g(\ln(2a))=2a-2a\ln(2a)-b
\]
当 \(a\geqslant\frac\e2\) 时，\(g'(x)\leqslant0\)，\(g\) 在 \([0,1]\) 上递减，最小值为
\[
g(1)=\e-2a-b
\]
综上
\[
\min_{x\in[0,1]}g(x)=
\begin{cases}
1-b,&a\leqslant\frac12,\\
2a-2a\ln(2a)-b,&\frac12<a<\frac\e2,\\
\e-2a-b,&a\geqslant\frac\e2.
\end{cases}
\]

\item 由 \(f(1)=0\)，得
\(\e-a-b-1=0\)，\(b=\e-a-1\).
代入 \(f(x)\) 并整理：
\[
f(x)=\e^x-1-(\e-1)x+a x(1-x)
\]
若 \(x\in(0,1)\) 是零点，则 \(x(1-x)>0\)，所以
\[
a=h(x):=\frac{1+(\e-1)x-\e^x}{x(1-x)}
\]
利用指数函数的幂级数展开，
\[
1+(\e-1)x-\e^x=\sum_{n=2}^{\infty}\frac{x-x^n}{n!}
=x(1-x)\sum_{n=2}^{\infty}\frac{1+x+\cdots+x^{n-2}}{n!}
\]
所以
\[
h(x)=\sum_{n=2}^{\infty}\frac{1+x+\cdots+x^{n-2}}{n!}
\]
当 \(0<x_1<x_2<1\) 时，右端各项均不减，且从 \(n=3\) 项起严格增加，所以 \(h(x_1)<h(x_2)\)，即 \(h\) 在 \((0,1)\) 上严格递增. 又
\(\lim_{x\to0^+}h(x)=\sum_{n=2}^{\infty}\frac1{n!}=\e-2\)，\(\lim_{x\to1^-}h(x)=\sum_{n=2}^{\infty}\frac{n-1}{n!}=1\).
由于 \(h\) 连续且严格递增，其值域为
\[
h((0,1))=(\e-2,1)
\]
因此若存在 \(x\in(0,1)\) 使 \(f(x)=0\)，必有 \(a\in(\e-2,1)\). 反过来，任意 \(a\in(\e-2,1)\) 都由 \(h\) 的连续性和值域性质对应某个 \(x\in(0,1)\)，使 \(a=h(x)\)，进而 \(f(x)=0\). 所以
\[
a\in(\e-2,1)
\]
\end{enumerate}
\end{solution}
